\documentclass[11pt]{article}
% Unique hyperref destinations for the many unlabelled [H] tables in this paper.
\PassOptionsToPackage{hypertexnames=false}{hyperref}
% Shared preamble for the hanzo-kernel Paper 07 series.
% One style, one vocabulary, inputted by every paper so the series reads as one voice.
\usepackage[margin=1in]{geometry}
\usepackage{booktabs}
\usepackage{amsmath}
\usepackage{amssymb}
\usepackage{graphicx}
\usepackage{xcolor}
\usepackage{hyperref}
\hypersetup{colorlinks=true, linkcolor=blue, urlcolor=blue, citecolor=blue}
\usepackage{enumitem}
\usepackage{listings}
\usepackage{array}

\definecolor{kw}{rgb}{0.13,0.29,0.53}
\definecolor{cmt}{rgb}{0.40,0.40,0.40}
\definecolor{str}{rgb}{0.16,0.50,0.16}
\lstset{
  basicstyle=\ttfamily\footnotesize,
  keywordstyle=\color{kw}\bfseries,
  commentstyle=\color{cmt}\itshape,
  stringstyle=\color{str},
  showstringspaces=false,
  breaklines=true,
  columns=fullflexible,
  frame=single,
  framesep=4pt,
}
\lstdefinelanguage{Rust}{
  morekeywords={fn,let,mut,pub,struct,impl,trait,enum,match,if,else,for,while,loop,return,use,mod,crate,self,super,async,await,where,type,const,static,unsafe,ref,move,dyn,as,true,false},
  morecomment=[l]{//}, morecomment=[s]{/*}{*/}, morestring=[b]", sensitive=true,
}

\newcommand{\x}{\ensuremath{\times}}
\newcommand{\dpa}{\texttt{dp4a}}
\newcommand{\hk}{\texttt{hanzo-kernel}}
\newcommand{\code}[1]{\texttt{#1}}

% Absorb minor prose overfulls without rivers; keep line-breaking sane.
\tolerance=800
\emergencystretch=3em


\usepackage{amsthm}
\usepackage{float}
\theoremstyle{definition}
\newtheorem{finding}{Finding}

\newcommand{\rhobar}{\ensuremath{\bar{\rho}}}
\newcommand{\mfjp}{\textsc{Mfjp}}

\title{\textbf{Does the Mean-Field Machinery Earn Its Keep?}\\[2mm]
\large An ablation of calibration and reputation in a production LLM judge panel}
\author{Hanzo AI --- ML Systems}
\date{}

\begin{document}
\maketitle

\begin{abstract}
Hanzo ships a Mean-Field Judge Panel (\mfjp): a population of diverse LLM judges whose raw
scores are per-judge z-score calibrated and then combined with multiplicative-weights
reputation. We ask whether the two mean-field mechanisms --- calibration and reputation ---
earn their keep against the obvious alternatives, and we set out to falsify rather than
confirm. Two tracks: 164 HumanEval candidates labelled by \emph{executing} their tests,
scored by three real judges from different model families under the verbatim production
rubric; and a Monte-Carlo study ($400$ replicates) that drives the \emph{shipped}
\texttt{panelScore} through regimes chosen so the panel can lose. The ablation is surgical:
identical online calibration with reputation weights pinned to $1$ isolates reputation, and a
plain mean of raw scores isolates calibration.
\textbf{The headline is negative, and it holds on real data.} On the real judges, a plain mean
of the same three scores significantly outperforms the shipped panel (AUC $0.859$ vs $0.824$,
paired $\Delta=+0.035$, 95\% CI $[+0.014,+0.057]$), while the calibration-only ablation is
indistinguishable from \mfjp{} --- so reputation contributes nothing and calibration actively
costs. Two measured properties explain it: the production judges differ in bias but barely in
\emph{scale} (sd ratio $1.19$), and their errors are strongly correlated ($\rhobar=0.729$)
despite spanning three model families. In simulation the same picture holds --- under the
panel's own assumptions a plain mean matches \mfjp{} to $0.0004$ Spearman and cost-matched
self-consistency beats it --- and we identify scale heterogeneity as the \emph{necessary
condition} for the machinery to pay: when we manufacture it, the ladder is clean
($0.938 \to 0.948 \to 0.958$), but a batch Gaussian EM estimator and the best single judge at
one third the cost both still beat \mfjp{}. At the correlation production actually exhibits, a
single judge at one third the cost is statistically indistinguishable from the panel
($\Delta$AUC $-0.004$, $[-0.055,+0.045]$).
Reputation is actively \emph{harmful} when the majority is systematically biased: it
down-weights the one accurate dissenter to $0.26$ while the two confounded judges keep
$0.41$--$0.46$, costing $0.039$ Spearman versus calibration alone. The panel's one clear,
large win is adversarial: when responses are engineered to game a single judge, \mfjp{} scores
$0.918$ against $0.553$ for cost-matched self-consistency --- diversity, not the mean-field
machinery, is what buys this. We also report that the claimed ``free confidence signal'' is
uninformative in benign simulation ($\rho = 0.036$) and \emph{negative} under distribution
drift, though it does work on real data ($+0.331$); and that z-scoring destroys absolute
calibration: at a realistic $0.9$ base rate \mfjp{} attains $0.657$ accuracy and $0.456$ ECE
where a plain mean attains $0.979$ and $0.190$. Finally, three defects in the shipped code: the live config names a retired model so
the ``3-judge'' panel silently runs as two; the process-global mutex is held across all
sequential judge HTTP calls, serialising judged traffic for up to $90$\,s; and calibration
state is in-process only, so every deploy re-enters cold start.
\end{abstract}

\section{The question}

A single LLM judge has one model's bias, one point of failure, and one thing to game. The
shipped answer at Hanzo is the Mean-Field Judge Panel, live as
\texttt{[claude-3-5-haiku, gpt-4o-mini, llama-3.1-8b]} at \texttt{sample 0.10}. It does two
things beyond averaging a committee:

\begin{enumerate}[nosep]
\item \textbf{Calibration.} Each judge's raw score is z-scored against that judge's own
running mean and standard deviation (Welford, online) and logistic-squashed, so a harsh
judge's $0.6$ and a lenient judge's $0.8$ both read as ``above my own bar.''
\item \textbf{Reputation.} Each judge carries a reliability weight updated multiplicatively,
$w \leftarrow (1-\pi)\,w e^{-\lambda|c-r|} + \pi$, with $\lambda=2.0$, $\pi=0.1$,
$w_{\min}=0.05$. A judge that tracks the calibrated consensus keeps its weight; one that
drifts is exponentially suppressed.
\end{enumerate}

The reward is the reliability-weighted calibrated consensus; panel disagreement is returned as
a free confidence signal. The design is computationally cheap by construction --- four floats
per judge, $O(1)$ per call.

The question this paper answers is not ``does a panel beat one judge'' (it usually does, and a
companion paper measured that). It is narrower and more consequential:
\textbf{does the mean-field machinery earn its keep over simply averaging the same diverse
judges?} And at equal budget --- three judges cost three times one judge --- does the panel
beat what else that budget could buy?

\subsection{What the prior work established, and what it did not}

A companion Hanzo study measured judge panels against executable ground truth and found that a
panel buys \emph{variance}, not accuracy: a five-judge majority reached $0.750$ against a mean
single-judge $0.685$, but the best single rubric alone reached $0.787$. It also reported a
negative: a mean-field construction in which judges \emph{see} the aggregate and revise
\emph{herds}, scoring below plain majority.

Two things leave our question open. First, that study's five judges shared one model and
differed only by rubric, so it measured rubric diversity, not model diversity --- a limitation
it flagged itself. Second, and more importantly, its ``mean-field'' arm is a
\emph{different algorithm} from the shipped \mfjp{}: it exposes the consensus to the judges,
whereas \mfjp{} aggregates independent scores post hoc and never lets judges see each other.
The herding result therefore says nothing about the shipped code. Nobody has ablated it.

\subsection{The shipped test begs the question}

The repository already contains a Monte-Carlo test,
\texttt{TestMFJPRecoversGroundTruth}. It is confirmatory in three ways, and we flag this
because it is the existing evidence base. Its judges have \emph{independent} Gaussian errors
around the truth --- which is precisely the panel's own generative assumption. Its single
adversary emits \emph{pure noise}, the easiest possible case for reputation to suppress. Its
metric is Spearman rank correlation, which is invariant to exactly the per-judge additive bias
that calibration removes, so calibration cannot be penalised. It then \emph{asserts} that
\mfjp{} wins, comparing only against a naive mean that conflates calibration and reputation
into one arm. We do not assert. We measure, and we choose conditions that can falsify.

\section{Method}

\subsection{Provenance}

All results drive the shipped algorithm. \texttt{router\_judge\_panel.go} was copied verbatim
from the published module \texttt{github.com/hanzoai/ai@v1.826.4} --- checksum-verified from
the module zip, byte-identical below the \texttt{package} line --- into a standalone harness
whose only addition is a shim supplying the two external symbols the file references. The
algorithm is untouched; \texttt{panelScore} is the production function.

\subsection{Track A --- real judges, executed ground truth}

We take 164 HumanEval problems and the candidate program generated for each by
Qwen2.5-Coder-7B, with pass/fail labels obtained by \emph{executing} the evalplus test suite
($89$ pass, $75$ fail; base rate $0.543$, so a constant-\textsc{pass} predictor scores
$0.543$). Three judges from genuinely different model families score each candidate using the
\emph{verbatim} production rubric from \texttt{judgeRubric} --- same system prompt, same strict
JSON contract, \texttt{temperature 0}, \texttt{max\_tokens 256}. No judge sees the tests or
the labels. This closes the model-diversity gap the prior study flagged.

One substitution was forced on us and is itself a finding (\S\ref{sec:defects}): the live
config names \texttt{claude-3-5-haiku}, which the gateway no longer serves. We substituted its
successor \texttt{claude-haiku-4.5}.

\subsection{Track B --- simulation that can falsify}

Track A gives one point in the space. Only simulation can \emph{sweep} the inter-judge error
correlation \rhobar{} and locate the boundary where the panel stops paying. We generate

\[
\text{raw}_{ij} \;=\; \mathrm{clip}_{[0,1]}\!\left(q_i + b_j + g_j c_i + \varepsilon_{ij}\right),
\qquad \varepsilon_{ij}\sim\mathcal N(0,s_j^2),\quad c_i\sim\mathcal N(0,1),
\]

where $q_i$ is latent quality, $b_j$ is judge $j$'s bias, and $c_i$ is a \emph{shared
confound} --- a property such as verbosity or fluency that judges reward but that is
orthogonal to quality. The error $e_{ij}=b_j+g_jc_i+\varepsilon_{ij}$ is independent across
judges only when $g=0$. The pairwise error correlation is
$\rho_{jk} = g_jg_k/\sqrt{(g_j^2+s_j^2)(g_k^2+s_k^2)}$, so $g$ is the dial that moves us from
the panel's assumed world to the one real LLM judges inhabit, where shared pretraining and
shared RLHF make errors correlate. Dawid--Skene assumes conditional independence too, so it
should degrade alongside. We report the \emph{realized} \rhobar{} measured from the generated
data, not the nominal target, because clipping attenuates it.

$400$ independent replicates of $600$ items each. Online arms are scored on the post-warmup
half so that calibration and reputation have converged --- a choice that \emph{favours}
\mfjp{}; cold start is reported separately.

\subsection{The surgical ablation}

This is the core of the design. Three arms differ by exactly one mechanism each:

\begin{itemize}[nosep]
\item \textbf{raw mean} --- unweighted mean of raw scores. No calibration, no reputation.
\item \textbf{z-mean} --- \emph{identical} online calibration to \mfjp{}: the same
\texttt{calibrateScore} function, the same \texttt{judgeStat} Welford state, the same
score-then-observe ordering. Reputation weights pinned to $1$.
\item \mfjp{} --- the shipped \texttt{panelScore}.
\end{itemize}

Any \mfjp{}\,--\,z-mean gap is attributable to \emph{reputation alone}; any z-mean\,--\,raw-mean
gap to \emph{calibration alone}. The shipped test conflates both into a single ``naive mean''
comparison and so cannot attribute credit.

\subsection{Competing arms}

Beyond the ablation ladder: \textbf{majority vote} (each judge votes against its own running
mean); \textbf{Dawid--Skene}, the classical latent-confusion-matrix EM aggregator, run in
\emph{batch} over all items --- a deliberate advantage over the online arms;
\textbf{Gaussian EM}, the principled continuous competitor, which is the maximum-likelihood
estimator under \emph{exactly the panel's own assumed model} (per-judge additive bias plus
per-judge Gaussian noise, conditionally independent) and therefore performs bias removal plus
inverse-variance weighting --- \mfjp{} is best understood as a cheap online heuristic
approximating this batch optimum; \textbf{cascade} (cheap judge first, escalate to the panel
only near its decision boundary); and two \textbf{cost-matched} arms that spend the panel's
three-call budget on one judge --- \textbf{self-consistency} (one judge sampled three times,
averaged) and a \textbf{homogeneous panel} (three samples of one judge fed to \mfjp{}).

\subsection{Metrics}

Spearman correlation with latent quality (ranking); AUC; Brier score and expected calibration
error over 10 bins (the confidence claim); accuracy; cost in judge calls. For the confidence
signal we correlate \mfjp{}'s returned confidence against \emph{rank} error --- $|{\rm
rank}(\hat q)-{\rm rank}(q)|$ --- because \mfjp{}'s z-squashed output is not on the same scale
as $q$ and a raw $|\hat q - q|$ would be scale-confounded and unfair to it. Differences are
reported as \emph{paired} deltas versus \mfjp{} with 95\% bootstrap confidence intervals
(4000 resamples over replicates; 5000 over items for Track A).

\section{Results}

\subsection{Track A --- real judges on executed ground truth}
\label{sec:trackA}

Three real judges scored 164 execution-labelled candidates. $44$ of $492$ calls ($8.9\%$)
failed to yield a parseable score even after retries and are recorded as abstentions. All arms
--- including \mfjp{} --- were run on the same matrix with abstentions imputed at the judge's
own mean, so no arm sees more information than another.

\begin{table}[H]\centering
\begin{tabular}{lrrrrr}
\toprule
judge & mean & sd & AUC & best acc & abstentions \\
\midrule
\texttt{claude-haiku-4.5} & 0.487 & 0.323 & 0.800 & 0.744 & 18 \\
\texttt{gpt-4o-mini}      & 0.734 & 0.343 & 0.821 & 0.805 & 16 \\
\texttt{llama-3.1-8b}     & 0.763 & 0.386 & 0.750 & 0.750 & 10 \\
\bottomrule
\end{tabular}
\caption{Per-judge behaviour on 164 items, base rate $0.543$.}
\end{table}

Two measurements decide how the rest of this paper applies to production.

\begin{finding}
\label{f:scale}
\textbf{The production judges differ in bias but not in scale.} Their means span $0.487$ to
$0.763$ --- a large $0.28$ offset gap, with Haiku much harsher --- but their standard
deviations are $0.323$, $0.343$, $0.386$, a max/min ratio of only $\mathbf{1.19}$. R8
established that scale heterogeneity is the \emph{necessary condition} for calibration and
reputation to pay. It is absent here. The mean-field machinery is deployed in precisely the
regime where the simulation predicts it is inert.
\end{finding}

\begin{finding}
\label{f:rho}
\textbf{The real judges are heavily correlated: $\rhobar = 0.729$}, with pairwise score
correlations of $0.575$, $0.533$ and $0.655$. Model-family diversity does \emph{not} deliver
error independence --- shared pretraining data and shared preference tuning dominate. Reading
$\rhobar=0.73$ off the R2 sweep places production near the far end of the curve, where the
diversity premium has decayed from $0.044$ to $0.026$.
\end{finding}

\begin{table}[H]\centering
\begin{tabular}{lrrrrrr}
\toprule
arm & AUC & best acc & acc@0.5 & Brier & ECE & cost \\
\midrule
\mfjp{} (calib+reputation)   & 0.824 & 0.774 & 0.738 & 0.183 & 0.135 & 3.00 \\
z-mean diverse (calib only)  & 0.827 & 0.780 & 0.738 & 0.181 & 0.131 & 3.00 \\
\textbf{raw mean diverse}    & \textbf{0.859} & \textbf{0.817} & 0.750 & \textbf{0.166} & 0.121 & 3.00 \\
\textbf{Gaussian EM (batch)} & \textbf{0.863} & \textbf{0.823} & \textbf{0.799} & 0.177 & 0.140 & 3.00 \\
Dawid--Skene (batch EM)      & 0.736 & 0.732 & 0.457 & 0.542 & 0.542 & 3.00 \\
majority vote                & 0.821 & 0.750 & 0.750 & 0.175 & \textbf{0.096} & 3.00 \\
cascade                      & 0.806 & 0.756 & 0.738 & 0.187 & 0.132 & \textbf{1.50} \\
single: \texttt{claude-haiku-4.5} & 0.800 & 0.744 & 0.732 & 0.190 & 0.123 & 1.00 \\
single: \texttt{gpt-4o-mini} & 0.821 & 0.805 & 0.652 & 0.201 & 0.204 & 1.00 \\
single: \texttt{llama-3.1-8b}& 0.750 & 0.750 & 0.701 & 0.257 & 0.251 & 1.00 \\
\bottomrule
\end{tabular}
\caption{Real judges, executed ground truth, $n=164$, $\rhobar=0.729$.}
\end{table}

Paired bootstrap over items ($5000$ resamples), $\Delta$AUC versus \mfjp{}:
raw mean $+0.0348$ $[+0.0142,+0.0567]^{*}$;
Gaussian EM $+0.0381$ $[+0.0072,+0.0699]^{*}$;
z-mean $+0.0022$ $[-0.0054,+0.0099]$;
majority $-0.0039$ $[-0.0323,+0.0240]$;
\texttt{gpt-4o-mini} alone $-0.0039$ $[-0.0551,+0.0447]$;
\texttt{claude-haiku-4.5} alone $-0.0247$ $[-0.0693,+0.0207]$;
Dawid--Skene $-0.0884$ $[-0.1401,-0.0356]^{*}$.

\begin{finding}
\textbf{On real judges and real executed ground truth, a plain mean of the same three judges
significantly outperforms the shipped panel}: AUC $0.859$ versus $0.824$, $\Delta=+0.0348$
with the interval excluding zero. The z-mean ablation is statistically indistinguishable from
\mfjp{} ($+0.0022$), so reputation contributes nothing; the raw-mean-versus-z-mean gap shows
the loss is attributable to \emph{calibration}, which --- consistent with
Finding~\ref{f:scale} --- has no scale heterogeneity to correct and merely discards
information.
\end{finding}

\begin{finding}
\textbf{At $\rhobar=0.73$, a single judge at one third the cost is statistically
indistinguishable from the panel.} \texttt{gpt-4o-mini} alone scores AUC $0.821$ against the
panel's $0.824$ ($\Delta=-0.0039$, $[-0.0551,+0.0447]$), and on best-threshold accuracy it is
higher ($0.805$ vs $0.774$). This is the cost-matched boundary the simulation predicted: the
diversity premium does not survive the error correlation that real judges actually exhibit.
\end{finding}

\paragraph{Robustness.} Dropping every item on which any judge abstained leaves $122$
complete cases and no imputation anywhere. The conclusions are unchanged: raw mean
$+0.0280$ $[+0.0072,+0.0503]^{*}$, Gaussian EM $+0.0415$ $[+0.0104,+0.0767]^{*}$, z-mean
$+0.0016$ $[-0.0074,+0.0110]$ versus \mfjp{}. The result is not an artifact of how
abstentions were handled.

\paragraph{The confidence signal does work here.} On real data \mfjp{}'s disagreement signal
correlates $+0.331$ with correctness: mean absolute error is $0.344$ on the high-confidence
half against $0.448$ on the low-confidence half. This is the one component that survives Track
A intact, and it is consistent with the simulation --- $\rhobar=0.73$ with a $0.28$ spread in
judge means is exactly the structural disagreement the signal is able to detect.

\paragraph{Dawid--Skene fails on real data} (AUC $0.736$, Brier $0.542$). Binarising a
continuous judge score at each judge's median discards the magnitude information that carries
most of the signal here. The classical aggregator is the wrong tool when judges emit graded
scores rather than categorical labels --- a point worth recording, since it is the aggregator
the literature would nominate.

\subsection{R1 --- the panel's own assumptions}

Independent errors, homogeneous scales, realized $\rhobar = 0.047$: the regime the design
assumes and the shipped test simulates.

\begin{table}[H]\centering
\begin{tabular}{lrrrrrr}
\toprule
arm & Spearman & AUC & Brier & ECE & acc & cost \\
\midrule
\mfjp{} (calib+reputation)      & 0.970 & 0.990 & 0.112 & 0.258 & 0.943 & 3.00 \\
z-mean diverse (calib only)     & 0.971 & 0.990 & 0.112 & 0.258 & 0.943 & 3.00 \\
raw mean diverse                & 0.971 & 0.990 & 0.092 & 0.209 & 0.943 & 3.00 \\
Gaussian EM (batch)             & 0.933 & 0.972 & 0.099 & 0.163 & 0.906 & 3.00 \\
Dawid--Skene (batch EM)         & 0.901 & 0.976 & 0.059 & 0.051 & 0.934 & 3.00 \\
majority vote                   & 0.903 & 0.976 & 0.052 & 0.030 & 0.935 & 3.00 \\
cascade                         & 0.947 & 0.987 & 0.106 & 0.243 & 0.943 & 1.74 \\
\textbf{self-consistency ($\times$3)} & \textbf{0.973} & \textbf{0.991} & 0.088 & 0.197 & 0.945 & 3.00 \\
homogeneous panel ($\times$3)   & 0.972 & 0.990 & 0.113 & 0.260 & 0.944 & 3.00 \\
single judge (best, oracle)     & 0.932 & 0.973 & 0.115 & 0.172 & 0.846 & 1.00 \\
single judge (random draw)      & 0.926 & 0.972 & 0.121 & 0.176 & 0.829 & 1.00 \\
\bottomrule
\end{tabular}
\caption{R1, independent errors, $\rhobar=0.047$, 400 replicates.}
\end{table}

\begin{finding}
Under its own assumptions the mean-field machinery adds \textbf{nothing}. Paired deltas versus
\mfjp{}: z-mean $+0.0002$ $[+0.0001,+0.0002]$, raw mean $+0.0004$ $[+0.0004,+0.0005]$. Both
intervals exclude zero --- so the plain mean is \emph{significantly better} --- but the effect
size is four ten-thousandths of a correlation unit. The honest statement is that they are
indistinguishable in practice, and that neither calibration nor reputation buys anything here.
\end{finding}

The reason is structural rather than empirical. Constant per-judge bias shifts every item
equally, so it cannot affect \emph{ranking}; Spearman is invariant to it. Calibration's stated
purpose --- removing each judge's standard --- is therefore irrelevant to rank quality when
judges differ only in offset. That leaves absolute level, which z-scoring actively destroys
(\S\ref{sec:abs}).

\begin{finding}
At equal cost the panel loses. Self-consistency --- one judge sampled three times --- scores
$0.973$ versus $0.970$ ($+0.0023$, $[+0.0020,+0.0026]$). Three calls spent on one judge beat
three calls spread across three judges in benign conditions.
\end{finding}

\subsection{R2 --- sweeping \rhobar{}: where the premium goes}

\begin{table}[H]\centering
\begin{tabular}{rrrrrr}
\toprule
\rhobar{} & \mfjp{} & z-mean & Gaussian EM & 1 judge & self-cons. \\
\midrule
0.047 & 0.970 & 0.971 & 0.933 & 0.926 & 0.973 \\
0.124 & 0.962 & 0.963 & 0.925 & 0.919 & 0.965 \\
0.203 & 0.953 & 0.953 & 0.916 & 0.911 & 0.955 \\
0.284 & 0.941 & 0.941 & 0.904 & 0.900 & 0.943 \\
0.367 & 0.926 & 0.926 & 0.890 & 0.887 & 0.927 \\
0.452 & 0.906 & 0.906 & 0.871 & 0.868 & 0.907 \\
0.540 & 0.877 & 0.878 & 0.845 & 0.842 & 0.878 \\
0.632 & 0.835 & 0.835 & 0.806 & 0.803 & 0.836 \\
0.728 & 0.764 & 0.764 & 0.740 & 0.738 & 0.764 \\
0.831 & 0.621 & 0.621 & 0.607 & 0.602 & 0.621 \\
\bottomrule
\end{tabular}
\caption{Spearman versus realized inter-judge error correlation. 400 replicates per row.}
\end{table}

Two readings. First, the \textbf{diversity premium} --- \mfjp{} minus a randomly drawn single
judge --- decays monotonically from $0.044$ at $\rhobar=0.05$ to $0.019$ at $\rhobar=0.83$: as
judges correlate, the third judge stops adding information, exactly as theory predicts.
Second, and more damaging, \mfjp{} and z-mean are \textbf{identical to three decimal places at
every \rhobar{}}, and cost-matched self-consistency equals or beats \mfjp{} at every
\rhobar{}. The reputation mechanism contributes nothing anywhere along this sweep.

\subsection{R3 --- reputation is harmful when the majority is biased}

Two judges load on the shared confound ($g=0.35$); one is accurate ($g=0$). Realized
$\rhobar=0.315$.

\begin{table}[H]\centering
\begin{tabular}{lrrrrrr}
\toprule
arm & Spearman & AUC & Brier & ECE & acc & cost \\
\midrule
\mfjp{}                        & 0.768 & 0.890 & 0.155 & 0.131 & 0.794 & 3.00 \\
\textbf{z-mean (calib only)}   & \textbf{0.807} & 0.910 & 0.149 & 0.151 & 0.814 & 3.00 \\
raw mean                       & 0.790 & 0.900 & 0.132 & 0.069 & 0.809 & 3.00 \\
Gaussian EM (batch)            & 0.629 & 0.818 & 0.182 & 0.096 & 0.737 & 3.00 \\
Dawid--Skene (batch EM)        & 0.727 & 0.875 & 0.220 & 0.222 & 0.775 & 3.00 \\
majority vote                  & 0.756 & 0.891 & 0.134 & 0.111 & 0.777 & 3.00 \\
self-consistency ($\times$3)   & 0.632 & 0.820 & 0.178 & 0.088 & 0.742 & 3.00 \\
\textbf{single judge (best)}   & \textbf{0.928} & \textbf{0.972} & 0.096 & 0.149 & \textbf{0.904} & \textbf{1.00} \\
single judge (random)          & 0.720 & 0.865 & 0.155 & 0.116 & 0.791 & 1.00 \\
\bottomrule
\end{tabular}
\caption{R3, biased majority with one accurate dissenter, $\rhobar=0.315$, 400 replicates.}
\end{table}

\begin{finding}
Reputation \textbf{actively hurts}: \mfjp{} $0.768$ against z-mean $0.807$, a loss of $0.039$
attributable to the reputation mechanism alone. The mechanism is visible in the weights. After
600 items the accurate dissenter holds weight $\mathbf{0.260}$ while the two confounded judges
hold $\mathbf{0.412}$ and $\mathbf{0.460}$. Reputation cannot distinguish ``drifting from the
consensus because it is wrong'' from ``drifting from the consensus because the consensus is
wrong,'' so it suppresses the only judge telling the truth by a factor of $1.7$.
\end{finding}

Note also that the best single judge --- the honest one --- scores $0.928$ at \emph{one third}
the cost, beating every three-judge arm. When a shared bias exists and one judge is free of
it, aggregation is the wrong move entirely; judge selection is the right one.

\subsection{R4 --- adversarial gaming: the panel's real win}

$12\%$ of items are engineered to spike one judge while true quality is low.

\begin{table}[H]\centering
\begin{tabular}{lrrrr}
\toprule
arm & Spearman & AUC & acc & cost \\
\midrule
\mfjp{}                       & \textbf{0.918} & \textbf{0.990} & 0.946 & 3.00 \\
z-mean (calib only)           & 0.910 & 0.989 & 0.945 & 3.00 \\
raw mean                      & 0.902 & 0.987 & 0.944 & 3.00 \\
Gaussian EM (batch)           & 0.928 & 0.977 & 0.915 & 3.00 \\
Dawid--Skene (batch EM)       & 0.868 & 0.979 & 0.930 & 3.00 \\
majority vote                 & 0.841 & 0.974 & 0.941 & 3.00 \\
\textbf{self-consistency ($\times$3)} & \textbf{0.553} & \textbf{0.797} & 0.833 & 3.00 \\
homogeneous panel ($\times$3) & 0.551 & 0.796 & 0.819 & 3.00 \\
single judge (random)         & 0.793 & 0.914 & 0.809 & 1.00 \\
\bottomrule
\end{tabular}
\caption{R4, reward hacking against one judge, $\rhobar=0.101$, 400 replicates.}
\end{table}

\begin{finding}
This is the panel's genuine, large win, and the only regime where reputation clearly pays.
\mfjp{} $0.918$ versus cost-matched self-consistency $0.553$ --- a gap of $0.365$
$[-0.371,-0.360]$. Sampling one judge three times is fooled every time, because it is the same
judge with the same exploitable bias. Reputation also earns its keep here: \mfjp{} beats
z-mean by $+0.0089$ and raw mean by $+0.0168$, both intervals excluding zero.
\end{finding}

The credit belongs mostly to \emph{diversity}, not to the mean-field machinery: raw mean over
diverse judges already reaches $0.902$ of the $0.918$. But the machinery does add a real
increment here, and this is the argument that justifies the panel's existence.

\subsection{R5 --- z-scoring destroys absolute calibration}
\label{sec:abs}

Quality $\sim\mathrm{Beta}(5,1)$: most responses are genuinely good, which is the regime a
production serving path actually occupies.

\begin{table}[H]\centering
\begin{tabular}{lrrrrr}
\toprule
arm & Spearman & AUC & Brier & ECE & acc \\
\midrule
\mfjp{}                       & 0.830 & 0.991 & 0.240 & 0.456 & 0.657 \\
z-mean (calib only)           & 0.831 & 0.991 & 0.240 & 0.456 & 0.655 \\
\textbf{raw mean}             & 0.828 & 0.992 & \textbf{0.056} & \textbf{0.190} & \textbf{0.979} \\
Gaussian EM (batch)           & 0.762 & 0.972 & 0.055 & 0.184 & 0.970 \\
Dawid--Skene (batch EM)       & 0.708 & 0.840 & 0.464 & 0.571 & 0.371 \\
\textbf{self-consistency}     & 0.845 & 0.992 & \textbf{0.048} & \textbf{0.159} & \textbf{0.981} \\
single judge (random)         & 0.691 & 0.972 & 0.080 & 0.182 & 0.907 \\
\bottomrule
\end{tabular}
\caption{R5, base rate $\approx 0.9$, $\rhobar=0.031$, 400 replicates.}
\end{table}

\begin{finding}
\mfjp{}'s reward is \textbf{purely relative} and its absolute calibration is catastrophic
where it matters most. Because \texttt{calibrateScore} z-scores against the judge's own running
mean and logistic-squashes, the output is forced toward a spread around $0.5$ regardless of how
good the responses actually are. When $90\%$ of responses are good, \mfjp{} attains accuracy
$0.657$ and ECE $0.456$; a plain mean of the same judges attains $0.979$ and $0.190$. Note
that AUC is \emph{identical} ($0.991$ vs $0.992$) --- ranking is preserved, absolute level is
destroyed. A reward consumer that thresholds the score, rather than only ranking by it, is
being misled on roughly one third of good responses.
\end{finding}

This is not a tuning problem; it is what z-scoring means. Any downstream use of the \mfjp{}
scalar as a probability or as an absolute quality gate is unsound by construction.

\subsection{R8 --- the one regime where the machinery earns its keep}

To avoid rigging the ablation, we constructed the case calibration is \emph{for}: judges that
differ in \emph{spread}, not merely offset (noise $s = 0.30$, $0.05$, $0.12$). A raw mean is
dominated by whichever judge swings widest; calibration equalises influence.

\begin{table}[H]\centering
\begin{tabular}{lrrrrr}
\toprule
arm & Spearman & AUC & Brier & ECE & cost \\
\midrule
\mfjp{}                        & 0.958 & 0.987 & 0.117 & 0.254 & 3.00 \\
z-mean (calib only)            & 0.948 & 0.983 & 0.120 & 0.246 & 3.00 \\
raw mean                       & 0.938 & 0.977 & 0.100 & 0.183 & 3.00 \\
\textbf{Gaussian EM (batch)}   & \textbf{0.983} & \textbf{0.995} & 0.089 & 0.223 & 3.00 \\
Dawid--Skene (batch EM)        & 0.883 & 0.982 & 0.045 & 0.041 & 3.00 \\
self-consistency ($\times$3)   & 0.860 & 0.940 & 0.117 & 0.130 & 3.00 \\
\textbf{single judge (best)}   & \textbf{0.983} & \textbf{0.995} & 0.116 & 0.228 & \textbf{1.00} \\
\bottomrule
\end{tabular}
\caption{R8, heterogeneous judge scales, $\rhobar=0.089$, 400 replicates.}
\end{table}

\begin{finding}
Both mechanisms do earn their keep here, and the ladder is clean and monotone:
raw mean $0.938 \to$ z-mean $0.948$ (calibration, $+0.0099$) $\to$ \mfjp{} $0.958$
(reputation, $+0.0099$), both intervals excluding zero. \textbf{Scale heterogeneity is the
necessary condition for the mean-field machinery to pay.} But even here \mfjp{} is beaten by
Gaussian EM ($0.983$, $+0.0259$) and matched by the best single judge at one third the cost.
\end{finding}

This makes the design decision empirical rather than theoretical: the machinery pays only if
the deployed judges actually differ in scale. \S\ref{sec:trackA} measures whether they do.

\subsection{The confidence signal}

Panel disagreement is advertised as a free confidence signal. Correlating it against rank
error across regimes:

\begin{table}[H]\centering
\begin{tabular}{lr}
\toprule
regime & $\rho(\text{confidence},\,-|\text{rank error}|)$ \\
\midrule
R1 independent errors            & $+0.036$ \\
R3 biased majority               & $+0.561$ \\
R4 reward hacking                & $+0.256$ \\
R5 high base rate                & $+0.034$ \\
R6 non-stationary quality        & $\mathbf{-0.143}$ \\
R8 heterogeneous scales          & $+0.271$ \\
\bottomrule
\end{tabular}
\caption{Informativeness of the returned confidence.}
\end{table}

\begin{finding}
The confidence signal is \textbf{uninformative in benign conditions} ($+0.036$) and
\textbf{actively misleading under drift} ($-0.143$: the panel is most confident when it is
most wrong). It is genuinely informative only when there is structural disagreement to detect
--- a biased faction ($+0.561$), an adversarial item ($+0.256$), or scale heterogeneity
($+0.271$). It should be used as a \emph{disagreement detector}, which is what it measures,
and never as a general-purpose reliability estimate.
\end{finding}

\subsection{Cold start}

Scored on the first $8\%$ of the stream with no warmup, \mfjp{} falls to $0.965$ while the
stateless raw mean holds $0.971$ and self-consistency $0.973$. The online state is a small
liability before it converges, which matters more than it appears (\S\ref{sec:defects}).

\section{Three defects in the shipped code}
\label{sec:defects}

These are independent of the statistics and were found by reading and exercising the shipped
artifact.

\paragraph{1. The panel silently runs with two judges, not three.}
The live configuration names \texttt{claude-3-5-haiku}. That model was retired on 2026-02-19
and the gateway returns HTTP 400 for it. In \texttt{panelScore} a judge whose call fails is
skipped by \texttt{continue} --- ``an erroring judge abstains; the panel stays resilient.''
The resilience is real, and it is precisely what conceals the fault: the advertised three-judge
mean-field panel has been running as a two-judge panel with no error surfaced. Judge
abstentions need a metric.

\paragraph{2. A process-global mutex is held across all judge network calls.}
\texttt{panelScore} takes \texttt{panelState.mu.Lock()} and \texttt{defer}s the unlock
\emph{before} the loop that issues the judge HTTP requests. The judges are called
\textbf{sequentially} --- there are no goroutines in the function --- and \texttt{judgeClient}
carries a $30$\,s timeout. So a single judged turn can hold a process-wide lock for up to
$3\times30 = 90$\,s, during which every other judged turn in the process blocks. Panel latency
is additive ($\sum$, not $\max$) and the panel is a global serialisation point. The fix is to
issue the calls concurrently outside the lock and take it only to fold results into state.

\paragraph{3. Calibration state is in-process and unpersisted, so every deploy resets it.}
\texttt{panelState} is a package-level variable with no persistence. Every pod restart returns
all judges to $n{=}0$, where \texttt{std()} falls back to a $0.2$ prior and the first score is
z-scored against a mean of zero --- a raw $0.7$ calibrates to $0.97$. At the production
\texttt{sample 0.10} rate, only one turn in ten contributes state, so the panel spends a long
window after every rollout in the cold-start regime of \S3.8, emitting inflated rewards.

\section{Recommendation}

The evidence supports one recommendation, in three parts.

\textbf{Keep the diverse panel; drop the mean-field machinery.} Replace the
reliability-weighted calibrated consensus with a plain mean of raw diverse judge scores. On
real judges against executed ground truth it beats \mfjp{} by $0.035$ AUC with the interval
excluding zero; in simulation it matches \mfjp{} to $0.0004$ Spearman under benign conditions,
beats it by $0.039$ when the majority is biased, and is dramatically better calibrated in
absolute terms ($0.056$ vs $0.240$ Brier, $0.979$ vs $0.657$ accuracy at a realistic base
rate). The necessary condition for the machinery to pay --- judges differing in scale --- is
measurably absent from the deployed panel (sd ratio $1.19$). Diversity is doing the
work that matters --- the entire adversarial robustness result survives, since raw mean over
diverse judges reaches $0.902$ of \mfjp{}'s $0.918$ in R4. Calibration and reputation are
worth reinstating only if the deployed judges are measured to differ in \emph{scale}, and
reputation only if adversarial gaming is an active threat; if it is, weight by
\emph{held-out agreement with ground truth}, never by agreement with the current consensus,
which is what makes R3 fail.

\textbf{Do not use the score as a probability.} The z-scored output is a ranking signal.
Any consumer thresholding it as an absolute quality gate should consume a raw mean instead.

\textbf{Fix the three defects} in \S\ref{sec:defects} regardless of the aggregation choice;
the concurrency one is a live production risk independent of this study.

If the goal is peak accuracy rather than architectural simplicity, the batch Gaussian EM
estimator dominated \mfjp{} in every regime where the machinery mattered at all
($+0.026$ in R8) and is the principled version of what \mfjp{} approximates --- \mfjp{}'s
defensible claim was never accuracy but that it is an $O(1)$-state streaming approximation to
it. That claim is true; the measurements show the approximation is not buying enough over a
plain mean to justify its failure modes.

\section{Threats to validity}

Track B is \emph{simulation} and its conclusions are conditional on the generative model. We
chose that model to be able to falsify the panel --- correlated errors, a systematically biased
majority, non-stationarity --- rather than to reproduce its assumptions, and we report the
regime (R8) where the machinery wins. But a real judge population may be structured in ways
this parameterisation does not capture; the R2 sweep is a statement about a one-parameter
confound family, not about all forms of dependence. Latent quality is a scalar, which
understates multi-dimensional disagreement.

Track A is $n=164$ single-seed with three judges in one domain (code) where an oracle exists;
the extrapolation to non-verifiable work is exactly the extrapolation the method cannot test.
Its candidates come from one generator, so failure modes are not representative of all model
errors. Substituting \texttt{claude-haiku-4.5} for the retired \texttt{claude-3-5-haiku} means
Track A measures the panel as \emph{intended}, not the panel as \emph{currently running}.

Bootstrap intervals quantify sampling variability only; they do not cover model
misspecification, which is the dominant uncertainty in Track B. Finally, several deltas here
are statistically significant and practically negligible --- that is a property of $400$
replicates, and we have tried throughout to report effect sizes rather than lean on interval
exclusion.

\section{Conclusion}

The Mean-Field Judge Panel's population of diverse judges is worth keeping: it is what defeats
reward hacking, and by a large margin. The mean-field machinery layered on top of that
population is not. On real judges against executed ground truth, plainly averaging the same
three scores beats the shipped panel by $0.035$ AUC; the calibration-only ablation is
indistinguishable from the full mechanism, so reputation contributes nothing and calibration
costs. In simulation the mechanisms are worth four ten-thousandths of a correlation unit, in
the wrong direction, and are inert across the whole \rhobar{} sweep; reputation is
affirmatively harmful when the majority shares a bias, for the structural reason that
agreement with the consensus is not evidence when the consensus is what is in error; and the
calibration step destroys the absolute scale that a thresholding reward consumer depends on.
The machinery earns its keep in exactly one measured condition --- judges that differ in scale
--- and that condition is measurably absent from the deployed panel, whose judges differ in
bias but share a scale to within a factor of $1.19$ and whose errors correlate at
$\rhobar=0.729$ despite three distinct model families. Diverse judges, plainly averaged, is
the configuration the evidence supports.

\section*{Reproduction}

The harness drives the shipped \texttt{panelScore} copied verbatim from the checksum-verified
module \texttt{github.com/hanzoai/ai@v1.826.4}. It contains the generative regimes, the
surgical ablation, all competing arms, the paired bootstrap, and the Track A collector that
records the per-judge per-item score matrix --- an artifact the prior instrument computed in
memory and discarded.

\end{document}
